% =====================================================================
%  6. Resultados
%  Origem: células 37, 38 e 39 do caderno
%  NOTA: o label sec:empiricos é o alvo dos apontamentos deixados nas
%  subseções de método (m02, m07, m08, m09, m11).
% =====================================================================
\section{Resultados}
\label{sec:empiricos}

Tendo em vista que os melhores resultados foram obtidos com o léxico, com os dois métodos
apoiados no analisador sintático e com o modelo mascarado, os demais não foram incorporados
à cadeia. Os descritores --- índices de complexidade sintática, comprimento de dependência e
embeddings --- não respondem à pergunta da correção, uma vez que medem estilo, custo de
processamento e proximidade temática, e as frases que os separam no corpus são todas
corretas. Os escores globais, SLOR e PLL, exigem um termo de comparação de que a avaliação
de uma frase isolada não dispõe, e o SLOR degenera ainda fora do vocabulário do corpus de
contagens. A surpresa por token, por sua vez, foi excluída por medição direta: seu pico
recaiu sobre o verbo correto da frase de referência.

Aos quatro métodos selecionados incorporaram-se três ajustes identificados nas subseções
anteriores: a exigência de verbo finito na raiz, a comparação de gênero no particípio
passivo e o tratamento do numeral por valor, dado que o numeral não porta traço de número.

Seguindo as etapas lógicas da avaliação, da ortografia à semântica, obtém-se a ordem devida
da cadeia. A verificação ortográfica ocupa a primeira posição por constituir o único
veredito exato e porque uma palavra fora do léxico não permanece contida nessa etapa: ela
deforma a árvore de dependências e, com ela, as duas etapas que a leem. A verificação
estrutural antecede a semântica porque o método do verbo mascarado devolve \texttt{n/a}
quando a raiz não é um verbo; a existência de predicado constitui, portanto, pré-requisito
da etapa seguinte, e o teste da raiz a verifica a custo desprezível.

Cada etapa devolve um conjunto de achados, e o veredito resulta da união
%
\begin{equation}
\mathrm{Loc}(\sent) = E_1(\sent) \cup E_2(\sent) \cup E_3(\sent) \cup E_4(\sent)
\qquad
\mathrm{Loc}(\sent) \neq \varnothing \;\Longrightarrow\; \text{frase reprovada}
\end{equation}
%
O veredito decorre, portanto, da vacuidade de um conjunto: nenhum valor é somado, ponderado
ou confrontado com limiar agregado. A cadeia dispõe de um único parâmetro calibrável, o
$\theta$ da etapa semântica, fixado em 5,0 nats, e ele determina apenas a admissão daquele
achado, sem participar do veredito final. As três primeiras etapas decidem por pertinência a
um conjunto ou por igualdade de valores, sem parâmetro algum.

\subsection*{Percurso da cadeia sobre uma frase}

Aplica-se a seguir a cadeia, etapa a etapa, à frase \emph{Há tres dias o incêndio atingiu
uma áreas de mata natvia}, que reúne três desvios de naturezas distintas. Todas as quatro
etapas são executadas, ainda que a primeira produza achado.

A etapa inicial avalia a pertinência ao léxico, resguardados os nomes próprios e os
numerais:
%
\begin{equation}
E_1(\sent) = \bigl\{\, t \;:\; \ell(t) \notin \lexico,\ \ \mathrm{pos}(t) \neq \texttt{PROPN},\ \ \ell(t) \notin \mathbb{N} \,\bigr\}
\end{equation}
%
O predicado é avaliado sobre uma palavra por vez, contra o conjunto. Das onze palavras, dez
pertencem a $\lexico$ e apenas \texttt{natvia} não, donde $E_1(\sent) =
\{\texttt{natvia}\}$. Cabe registrar a segunda palavra: a forma \texttt{tres}, sem acento,
\textbf{pertence} ao léxico, com $f = 1{.}690$ ocorrências, uma vez que a lista de
frequências foi construída a partir de texto autêntico da web, em que a grafia sem acento
ocorre com frequência suficiente para ser incorporada. A etapa não produz achado nessa
posição, e não por defeito de implementação: trata-se do comportamento previsto pela
fórmula, em que $\lexico$ constitui a definição operacional de erro.

Para o achado detectado, a segunda parte do método computa os candidatos e resolve o
desempate por frequência:
%
\begin{equation}
C(\texttt{natvia}) = \{\texttt{nativa},\ \texttt{naevia},\ \texttt{natia},\ \texttt{latvia}\}
\qquad
\hat{w} = \arg\max_{c \in C} f(c) = \texttt{nativa} \quad (f = 912)
\end{equation}

A segunda etapa verifica a existência de predicado, já sob a exigência de finitude:
%
\begin{equation}
\begin{split}
\mathrm{pred}(\sent) \iff\;
& \bigl[\mathrm{pos}(\raiz) \in \{\texttt{VERB},\texttt{AUX}\} \ \wedge\ \mathrm{VerbForm}(\raiz) = \texttt{Fin}\bigr] \\
\vee\; & \bigl[\exists\, c \in \mathrm{apoio}(\raiz) : \mathrm{VerbForm}(c) = \texttt{Fin}\bigr]
\end{split}
\end{equation}
%
Não há percurso a executar, dado que se inspeciona um único nó. A raiz é \texttt{atingiu},
com $\mathrm{pos} = \texttt{VERB}$ e $\mathrm{VerbForm} = \texttt{Fin}$, e não há auxiliar
nem cópula. Satisfeita a primeira cláusula, a segunda não é avaliada, e $E_2(\sent) =
\varnothing$.

A terceira etapa percorre os arcos e desdobra cada um nos traços exigidos por $T(r)$. Dos
dez arcos da frase, apenas três são admitidos, pois nos demais $T(r) = \varnothing$: as
palavras \texttt{de}, \texttt{mata} e \texttt{áreas} vinculam-se por \texttt{case},
\texttt{nmod} e \texttt{obj}, e nenhuma dessas relações impõe concordância. A
Tabela~\ref{tab:percurso} apresenta as comparações efetuadas.

\begin{table}[H]
\centering
\dimtable{
\begin{tabular}{@{}llccl@{}}
\toprule
\textbf{Arco} $(t,r,h)$ & $\tau$ & $\tracosp{\tau}{t}$ & $\tracosp{\tau}{h}$ & \textbf{Resultado} \\
\midrule
\multirow{2}{*}{(\texttt{o}, det, \texttt{incêndio})}
  & Number & Sing & Sing & iguais \\
  & Gender & Masc & Masc & iguais \\
\midrule
\multirow{2}{*}{(\texttt{incêndio}, nsubj, \texttt{atingiu})}
  & Number & Sing & Sing & iguais \\
  & Person & $\ausente$ & 3 & comparação suprimida \\
\midrule
\multirow{2}{*}{(\texttt{uma}, det, \texttt{áreas})}
  & Number & \textbf{Sing} & \textbf{Plur} & $\in E_3(\sent)$ \\
  & Gender & Fem & Fem & iguais \\
\bottomrule
\end{tabular}}
\caption{As comparações da etapa gramatical sobre os três arcos admitidos.}
\label{tab:percurso}
\end{table}

A quarta linha exemplifica a guarda $\tracosp{\tau}{\cdot} \neq \ausente$: o substantivo
\texttt{incêndio} não porta traço de pessoa, ao passo que o verbo apresenta
$\texttt{Person}=3$, de modo que a comparação é suprimida. O resultado não é de igualdade,
mas de ausência de termos a comparar. Resta, portanto, $E_3(\sent) = \{(\texttt{uma},
\texttt{áreas}, \texttt{Number})\}$.

A quarta etapa inspecciona uma única posição, a da raiz, mascarando-a e comparando a
distribuição obtida com a palavra escrita:
%
\begin{equation}
m = \ln P\bigl(\hat{v} \mid \sent_{\setminus\raiz}\bigr) - \ln P\bigl(v_\raiz \mid \sent_{\setminus\raiz}\bigr)
\end{equation}
%
Mascarada a raiz, o modelo devolve como formas esperadas \texttt{[destruiu, atingiu,
invadiu, atinge, em]}, nas quais o verbo escrito ocupa a segunda colocação, com margem
$m = 0{,}95$. Como $m < \theta = 5{,}0$, tem-se $E_4(\sent) = \varnothing$. Manifesta-se
aqui o único limiar de toda a cadeia: as três etapas anteriores decidem por pertinência a um
conjunto ou por igualdade de valores, ao passo que esta compara um valor real a um ponto de
corte, e $\theta$ constitui escolha do analista.

Reunindo os quatro conjuntos,
%
\begin{equation}
\mathrm{Loc}(\sent) =
\underbrace{\{\texttt{natvia}\}}_{E_1}
\ \cup\ \underbrace{\varnothing}_{E_2}
\ \cup\ \underbrace{\{(\texttt{uma},\,\texttt{áreas})\}}_{E_3}
\ \cup\ \underbrace{\varnothing}_{E_4}
\ \neq\ \varnothing
\end{equation}
%
donde a frase é reprovada, com dois achados, cada qual identificando as palavras envolvidas e
a etapa de origem.

\subsection*{Propagação de erro entre etapas}

A Tabela~\ref{tab:laudo} consolida o laudo dessa frase ao lado do de uma variante que contém
apenas o desvio de concordância.

\begin{table}[H]
\centering
\dimtable{
\begin{tabular}{@{}lll@{}}
\toprule
\textbf{Etapa} & \emph{\ldots uma áreas de mata nativa.} & \emph{Há tres dias \ldots uma áreas \ldots natvia.} \\
\midrule
ortografia & ok \quad --- & \textbf{!!} \quad \texttt{natvia} $\to$ \texttt{nativa} \\
estrutura  & ok \quad predicado presente & ok \quad predicado presente $^\dagger$ \\
gramática  & \textbf{!!} \quad \texttt{uma}(Sing) × \texttt{áreas}(Plur)
           & \textbf{!!} \quad \texttt{uma}(Sing) × \texttt{áreas}(Plur) $^\dagger$ \\
semântica  & ok \quad margem 0,65 & ok \quad margem 0,95 \\
\bottomrule
\end{tabular}}
\caption{Laudo de duas frases. $^\dagger$ marcado como \emph{árvore sob suspeita}.}
\label{tab:laudo}
\end{table}

A decisão de não bloquear tem consequência observável. A palavra \texttt{natvia} não
permanece contida na etapa inicial: ela deforma a árvore que as etapas seguintes leem,
conforme a Tabela~\ref{tab:deformacao}.

\begin{table}[H]
\centering
\dimtable{
\begin{tabular}{@{}lll@{}}
\toprule
\textbf{Arco produzido} & & \textbf{Efeito observado} \\
\midrule
\texttt{(tres, obj, Há)} & & \texttt{tres} analisado como objeto de \texttt{Há} \\
\texttt{(dias, flat:name, tres)} & & \emph{tres dias} lido como nome próprio composto \\
\texttt{(natvia, dep, atingiu)} & & \texttt{dep} é o rótulo de relação não classificável \\
\bottomrule
\end{tabular}}
\caption{Arcos deformados pela palavra fora do léxico.}
\label{tab:deformacao}
\end{table}

Registra-se por isso a marcação \emph{árvore sob suspeita} nas duas etapas que leem a
árvore, marcação que não indica incorreção, mas menor confiabilidade. No caso examinado a
deformação não atingiu o arco relevante, e o achado de concordância é legítimo, o que
constitui, contudo, contingência posicional e não garantia.

Duas lacunas, de naturezas distintas, delimitam o alcance da cadeia. A primeira é a escolha
lexical: \emph{A três dias} atravessa as quatro etapas sem detecção, pois \texttt{a} consta
do léxico, é preposição e não estabelece arco de concordância, e a etapa semântica não
inspeciona essa posição. Trata-se exatamente do caso que as regras escritas à mão detectam,
o que torna sua integração a extensão mais imediata da cadeia. A segunda é o erro factual:
\emph{A área cresceu de 1.284 para 340 hectares} é português correto e falso apenas em
relação à base que originou o texto. Nenhum método de natureza linguística alcança esse
nível, uma vez que a verificação exige confrontar o texto com os campos que o geraram.
